\section{Introduction}

Large Language Models (LLMs) have demonstrated impressive reasoning capabilities across diverse tasks, including mathematical problem solving~\citep{pei2025mathfusion,didolkar2024metacognitive}, commonsense inference~\citep{toroghi2024verifiable}, and open-domain question answering~\citep{zhao2023divknowqa}. 
Despite their generalization ability, LLMs often struggle in domain-specific scenarios due to the lack of grounded external knowledge, resulting in factual hallucinations and high inference costs~\citep{huang2025survey,wang2024factuality}. To address these limitations, recent efforts have explored integrating domain knowledge into LLM reasoning. A promising direction to overcome these limitations is Knowledge Graph Question Answering (KGQA)~\citep{dammu2025dynamic,saxena2020improving,choi2023nutrea}, which integrates relational structures into the reasoning process to provide factual grounding and structural interpretability. By combining the expressiveness of natural language with the precision of knowledge graphs, KGQA offers a scalable solution to improve factual consistency, reasoning transparency, and answer reliability~\citep{liu2025dual,yao2025learning}. 

Existing KGQA approaches can be broadly categorized into two paradigms according to how reasoning paths are constructed: retrieve-then-reason methods and dynamic path generation strategies. In the former, candidate paths are extracted prior to answer prediction, typically using Graph Neural Networks (GNNs) \citep{ma2025debate,yao2025learning} or rule-based heuristics \citep{fang2024karpa}. However, such methods exhibit limited adaptability, as GNNs fail to incorporate question-specific semantics during inference, while heuristic rules remain inherently inflexible and cannot support dynamic refinement of reasoning~\citep{liu2025dual,yao2025learning}. In contrast, dynamic path generation strategies unify retrieval and reasoning by constructing paths dynamically during question processing. These methods either leverage LLMs to iteratively generate paths through in-context learning or Chain-of-Thought (CoT)~\citep{sui2024fidelis,li2024framework}, or employ guided search techniques such as MCTS, where paths are incrementally expanded with a path scorer~\citep{ma2025deliberation,shen2025reasoning}. While offering greater flexibility, such approaches suffer from substantial computational overhead due to repeated LLM calls and limited evaluation accuracy, as static scorers cannot capture the evolving semantics of reasoning paths~\citep{chang2024survey,lee2024zero}.

This paper investigates the design of an adaptive KGQA framework to address challenges of computational inefficiency and limited path evaluation accuracy in dynamic reasoning. However, developing such a framework raises three key challenges: (1) \textit{How to modularize reasoning to reduce LLM overuse in search?} Inefficiency in dynamic KGQA stems from repeatedly invoking LLMs for relation retrieval and reasoning in multi-hop path construction~\citep{shen2025reasoning,long2025enhancing}. Although methods such as CoT and MCTS enable flexible exploration, they bind LLMs to every decision step, causing high inference cost and poor scalability. The challenge is to design a modular framework that leverages LLMs efficiently, guiding search without direct involvement at each step. (2) \textit{How to accurately evaluate evolving reasoning paths?} As multi-hop paths are incrementally constructed, their semantics evolve with each added relation and context. Yet existing methods rely on static scoring or shallow similarity metrics that fail to capture these semantic shifts~\citep{xu2024llm,sui2024fidelis}. This highlights the challenge of designing a path evaluation model that adaptively captures fine-grained changes conditioned on both the question and evolving relation sequence. (3) \textit{How to train a reliable evaluation model with limited supervision?} Accurate path ranking requires a well-calibrated scorer, yet dynamic reasoning generates many incomplete paths with only a few valid ones. This results in imbalanced, noisy supervision, especially for multi-hop questions where successful trajectories are scarce. While reinforcement learning has been explored~\citep{ma2024coevolving,zhai2024fine}, it suffers from sparse rewards and unstable optimization. The key challenge is to construct reliable learning signals from limited supervision to support adaptive scorer training.

% This paper investigates the design of an adaptive KGQA framework to address the challenges of computational inefficiency and limited path evaluation accuracy in dynamic reasoning. However, developing such a framework presents several key challenges: (1) \textit{How to modularize reasoning to reduce LLM overuse during search?} A major source of computational inefficiency in dynamic KGQA lies in repeatedly invoking LLMs for both relation retrieval and reasoning during multi-hop path construction~\citep{shen2025reasoning,long2025enhancing}. While methods such as CoT and MCTS provide flexible exploration, they tightly couple LLMs with each decision step, resulting in high inference costs and limited scalability. The key challenge is to design a modular reasoning framework that utilizes LLMs efficiently, guiding the search process without requiring direct involvement in every reasoning step.
% (2) \textit{How to accurately evaluate evolving reasoning paths?} As multi-hop reasoning paths are incrementally constructed, their semantics evolve with each newly added relation and contextual information. However, existing methods typically rely on static scoring functions or shallow similarity metrics, which fail to capture the nuanced semantic shifts that occur throughout the reasoning process~\citep{xu2024llm,sui2024fidelis}. This raises a key challenge: how to design a path evaluation model that adaptively captures fine-grained semantic changes conditioned on both the question and the evolving relation sequence.
% (3) \textit{How to train a reliable path evaluation model with limited supervision?} Accurate path ranking in KGQA hinges on a well-calibrated evaluation model. However, dynamic reasoning methods typically produce a large number of incomplete or irrelevant paths, with only a small subset corresponding to valid reasoning trajectories. This results in highly imbalanced and noisy supervision, especially for multi-hop questions where successful paths are extremely sparse. Although reinforcement learning has been explored to mitigate this issue~\citep{ma2024coevolving,zhai2024fine}, it frequently suffers from sparse rewards and unstable optimization. Therefore, the key challenge is how to construct meaningful learning signals from limited or implicit supervision to enable adaptive training of the path scorer.

To address these challenges, we propose \textbf{D}ynamically \textbf{A}daptive \textbf{M}CTS-based \textbf{R}easoning (\method{}), an efficient framework that integrates LLM-guided MCTS with context-aware semantic modeling for accurate and efficient KGQA.
\method{} employs an MCTS backbone, where an LLM-based planner dynamically guides path expansion by selecting semantically relevant relations at each step, thereby reducing the search space and improving answer identification.
For path evaluation, we introduce a lightweight Transformer-based scorer that jointly encodes the question and relation sequence via cross-attention, enabling context-sensitive plausibility estimation and capturing evolving semantics during multi-hop reasoning.
To mitigate supervision scarcity, \method{} further incorporates a dynamic pseudo-path mechanism that continuously adapts the scorer during search: partial paths from MCTS rollouts are ranked by predicted plausibility and converted into pseudo-path supervision pairs, amplifying signals from promising trajectories while suppressing noise from suboptimal ones.
Our contributions are summarized as follows:

% To address these challenges, we propose \textbf{D}ynamically \textbf{A}daptive \textbf{M}CTS-based \textbf{R}easoning (\method{}), an efficient framework that combines LLM-guided MCTS with context-aware semantic modeling for accurate and efficient KGQA.
% \method{} employs an MCTS backbone, where an LLM-based planner dynamically guides path expansion by selecting semantically relevant relations at each step, thereby reducing redundant LLM calls and constraining the search space for scalable, efficient reasoning.
% For path evaluation, we introduce a lightweight Transformer-based scorer that jointly encodes the question and relation sequence via cross-attention, enabling context-sensitive plausibility estimation and capturing evolving semantics during multi-hop reasoning.
% To mitigate supervision scarcity, \method{} incorporates a dynamic pseudo-path mechanism that continuously adapts the scorer during search. Partial paths sampled from MCTS rollouts are ranked by predicted plausibility and converted into pseudo-path supervision pairs, amplifying learning signals from promising trajectories while suppressing noise from suboptimal ones.
% Extensive experiments on benchmark KGQA datasets demonstrate that \method{} significantly outperforms state-of-the-art baselines. Our contributions are summarized as follows:
\begin{itemize}[itemsep=2pt,topsep=0pt,parsep=0pt,leftmargin=*]
\item We study adaptive path reasoning in KGQA, where the key challenges lie in capturing the evolving semantics of multi-hop reasoning paths and ensuring computational efficiency during search, motivating the need for dynamic and context-aware reasoning strategies.
\item We propose \method{}, a novel framework that integrates MCTS with a dynamically adapted path evaluation model, enhancing evaluation accuracy while maintaining computational efficiency.
\item We conduct extensive experiments across multiple KGQA benchmarks, demonstrating that \method{} consistently outperforms state-of-the-art methods.
\end{itemize}